\section{Prologue}\label{prologue}

I found a paper \cite{Woodhouse2024}; frf, matrices, \(\partial\).
Prof. Woodhouse' publications might be more relevant that Rayleigh's book.

\url{https://scholar.google.com/citations?user=2jBVQo4AAAAJ&hl=en}

Also see \url{https://euphonics.org/}



This document collects Gemini responses to prompts discussing prerequisities.

My primary motivation for reading Theory of Sound is to learn the
approximation for the first fundamental mode of structures \cite{Hodges1985}.
The static deflection curve is the deformation of the structure under a uniform
load, e.g., the force of gravity. Rayleigh's Method for estimating the
fundamental frequency of a structure is to use the Rayleigh quotient
corresponding to the static deflection curve.

Another reason to reading Theory of Sound is that it defined the
\emph{lingua franca} for a large proposition of the work related to
acoustics and vibrations done in the early 20th century. Readers must
know Rayleigh's notation `by heart'. Moreover Rayleigh assumed that his
readers were familiar with the famous 17th and 19th century inventors
and their intriguing inventions.

Humans hear sounds in the range 20~Hz to 20,000~Hz. Decibels (dB) are a
measure of wave amplitude, the displacement of particles from their
resting position. Meters \(m\), seconds \(s\), A medium's density and
elasticity determine the speed of sound.

\begin{itemize}
\item
  air (at 20\(^\circ \)C): \(\sim \)343 \(m/s\)
\item
  water: \(\sim \)1,480 \(m/s\)
\item
  steel: \(\sim \)5,960 \(m/s\)
\end{itemize}

As far as I know, heat is the primary cause of variation in sound speed;
temperature is inversely proportional to density.

At a hard or soft boundary, or at the interface between a different
medium, sound waves change at a

\begin{itemize}
\item
  Reflection: Bouncing off a surface, creating echoes or reverberation.
\item
  Refraction: Bending as it passes between zones of different
  temperatures or materials.
\item
  Diffraction: Bending around obstacles or spreading through openings.
\item
  Absorption: Losing energy as heat when hitting soft, porous materials.
\end{itemize}

Diffraction, e.g., Borda's mouthpiece, creates interesting
singularities. The list of intriguing physical phenomena includes
impedance, dispersion, sub-harmonics, Helmholtz's resonator and
thermoacoustics.

\section{Nomenclature}\label{nomenclature}
One motivation for defining Rayleigh's notation is to understand the work of
the other Rayleigh.
\[(D_x u ) + {({D_x u})}^T =
\begin{bmatrix}
2e & c & b\\
c & 2f & a\\
b & a & 2g\\
\end{bmatrix}.\]
\[e=\frac{\partial u}{\partial x},\quad f=\frac{\partial v}{\partial y},\quad
g=\frac{\partial w}{\partial z}\]
\[a=\frac{\partial w}{\partial y}+\frac{\partial v}{\partial z},\quad
b=\frac{\partial u}{\partial z}+\frac{\partial w}{\partial x},\quad
c=\frac{\partial v}{\partial x}+\frac{\partial u}{\partial y}\]
Dilatation (Volumetric Strain): \(\theta =e+f+g\)

Stress (Internal Forces per Unit Area)

To represent the internal forces acting across planes within a vibrating
solid Rayleigh tracked the six components of the stress system using
capital letters:

Normal Stresses: Represented by \(P, Q, R\) (pulling or pushing
perpendicular to the \(x, y, z\) planes).

Tangential (Shear) Stresses: Represented by \(S, T, U\) (acting parallel
to the faces of an infinitesimal cube).

When writing the potential energy of a deformed solid (the Strain Energy
Function, \(W\)), he linked these stresses directly to his strain
components via Hooke's Law using the notation:

\(P=\frac{\partial W}{\partial e},\quad Q=\frac{\partial W}{\partial f},\quad R=\frac{\partial W}{\partial g}\)

\(S=\frac{\partial W}{\partial a},\quad T=\frac{\partial W}{\partial b},\quad U=\frac{\partial W}{\partial c}\)

Body/External Forces: When expressing the equations of motion for an
elastic solid body, Rayleigh denoted the external components of force
acting per unit volume as \(X, Y, Z\).

The Equations of Motion: \(\rho \) is the mass density of the material.
Using his combined notation for stress and force components, Rayleigh
wrote the foundational 3D elasto-dynamic equations as:
\[\rho \ddot{u}=\rho X+
\frac{\partial P}{\partial x}+
\frac{\partial U}{\partial y}+
\frac{\partial T}{\partial z}\]

\begin{longtable}[]{@{}ccc@{}}
\caption{Rayleigh elasticity notation and contemporary equivalents.}\label{tbl:nomenclature1}
\tabularnewline%
quantity & symbol & modern \\
\midrule\noalign{}
\endfirsthead%
\toprule\noalign{}
quantity & symbol & modern \\
\midrule\noalign{}
\endhead%
\endlastfoot{}% chktxt 1
displacement & \(u, v, w\) & \(u_1, u_2, u_3\) \\
Normal Strain & \(e, f, g\) &
\(\epsilon_{xx}, \epsilon_{yy}, \epsilon_{zz}\) \\
Shear Strain & \(a, b, c\) &
\(\gamma_{yz}, \gamma_{xz}, \gamma_{xy}\) \\
Volumetric Strain & \(\theta = \Delta \) &
\(\epsilon_{v}= \text{Tr}(\boldsymbol{\epsilon})\) \\
Normal Stress & \(P, Q, R\) &
\(\sigma_{xx}, \sigma_{yy}, \sigma_{zz}\) \\
Shear Stress & \(S, T, U\) & \(\tau_{yz}, \tau_{xz}, \tau_{xy}\) \\
External Body Force & \(X, Y, Z\) & \(\mathbf{f} = b_{i}\) \\
Acoustic Pressure & \(p\) & \(p\), \(\Delta p\) \\
\end{longtable}


\begin{longtable}[]{@{}cc@{}}
\caption{Notation}\label{tbl:notation}\tabularnewline%
symbol & quantity modern \\
\midrule\noalign{}
\endfirsthead%
\toprule\noalign{}
symbol & quantity modern \\
\midrule\noalign{}
\endhead%
\endlastfoot%
\(T\) & kinetic energy \\
\(V\) & Strain energy \\
\(F\) & dissipation function \\
\(\psi_k\) & generalized coordinate \(k\) \\
\(y\), \(w\) & lateral displacement \\
A & amplitude \(A \cos( p t - \epsilon)\) \\
p & circular frequency \\
\(\epsilon \) & phase angle \\
\(\rho \) & mass density \\
\(\omega \) & cross sectional area \\
E & Young's \\
I & inertia \\
\(\kappa \) & radius of gyration \\
\end{longtable}


When writing the Strain Energy Function (W)---which measures the
potential energy stored in an elastic body when it is deformed.

Rayleigh formulated it using his strain components

(e, f, g for normal strains and a, b, c for shear strains)

alongside m and n:

\(2W=(m+n)(e^{2}+f^{2}+g^{2})+2(m-n)(f g+g e+ef)+n(a^{2}+b^{2}+c^{2})\)

Or in terms of the volumetric dilatation \(\theta = e + f + g\),

\(2W=(m-n)\theta^{2}+2n(e^{2}+f^{2}+g^{2})+n(a^{2}+b^{2}+c^{2})\)

\begin{longtable}[]{@{}ccc@{}}
\caption{Rayleigh's notation (blanks are intentional).}\label{tbl:archaic}
\tabularnewline%
quantity & symbol & modern \\
\midrule\noalign{}
\endfirsthead%
\toprule\noalign{}
quantity & symbol & modern \\
\midrule\noalign{}
\endhead%
\endlastfoot{}%
                 &\(m\)& \(\lambda + \mu \)  \\
shear modulus    &\(n\)& \(\mu = G\) \\
bulk modulus     &     & \(K = \lambda + \frac{2}{3} \mu \) \\
Poisson's ratio  &     & \(\nu = \frac{\lambda }{2 (\lambda + \mu )} \) \\
\end{longtable}

\section{Vocabulary Advance
Organizer}\label{vocabulary-advance-organizer}

\begin{longtable}[]{@{}ccc@{}}
\caption{Inventors and their inventions by chapter.}\label{tbl:nomenclature3}\tabularnewline%
Chapter & & dates, invention \\
\midrule\noalign{}
\endfirsthead%
\toprule\noalign{}
Chapter & & dates, invention \\
\midrule\noalign{}
\endhead%
\endlastfoot%
1 & Regnault & 1810--1878 \\
1 & Cagniard de la Tour & 1819, siren \\
2 & Wheatstone & 1827, kaleidophone \\
3 & Scheibler & 1834, tonometer \\
6 & Mersenne & 1588--1648, sonometer \\
7 & Terquem & 1831--1887 Savart's son rauque \\
9 & Bourget & 1822--1887 \\
10 & Chladni, Mathieu & \\
10 & Rudolph Koenig & 1832--1901 \\
\end{longtable}

Henri Victor Regnault (1810--1878) was a highly precise French chemist
and physicist, famous for his meticulous empirical measurements of the
physical properties of gases.

\textbf{His Contribution}: Regnault conducted massive, definitive
experiments on the velocity of sound in the atmosphere and inside long
pipes (famously utilizing the newly constructed Paris water mains). He
discovered that the speed of sound is not perfectly constant but varies
slightly depending on the intensity (loudness) of the sound wave and the
diameter of the tube it travels through.

\textbf{In Theory of Sound}: Rayleigh cites Regnault's measurements to
bridge the gap between idealized mathematical wave equations and the
messy reality of friction, thermal conduction, and atmospheric pressure
in real-world pipes.

\textbf{The Siren of Baron Charles Cagniard de la Tour} Invented by
French physicist Baron Charles Cagniard de la Tour in 1819, this device
was the first mechanical instrument capable of measuring the exact
absolute frequency (pitch) of a sound.

\textbf{How it worked}: It consisted of a pressurized air chamber with a
fixed metal disk (stator) on top, overlaid by a rotating metal disk
(rotor). Both disks were punched with a series of regularly spaced
holes.

\textbf{The Physics}: Crucially, the holes were drilled at opposing,
slanted angles. When air was forced upward from a bellows, it rushed
through the holes and physically pushed the top disk, making it spin. As
the disk spun, it rapidly cut off and opened the air stream, releasing
regular, rapid \textit{puffs} of air that the human ear perceives as a musical
tone.

\textbf{Rayleigh's Use}: Gears on the spinning disk connected to a dial
mechanism that precisely counted the number of rotations per second. By
multiplying the rotations by the number of holes, scientists could
calculate the exact frequency of the tone mathematically, providing an
invaluable tool for calibrating early acoustics experiments.

\textbf{The Mythological Name}:
Cagniard de la Tour named the device a siren because it could produce
sound underwater.

\textbf{Kaleidophone} Invented by Sir Charles Wheatstone in 1827, the
kaleidophone (literally \textit{beautiful sound-vessel}) was a device designed
to visually demonstrate the paths of vibrating objects using the
phenomenon of persistence of vision.

\textbf{How it worked}: It consisted of a thin, elastic metal rod
clamped firmly at one end. At the free upper end, a small, highly
polished or silvered glass bead was attached to reflect a point of
light.

\textbf{Rayleigh's Context}: When the rod was bowed or struck, it
vibrated transversely in two directions. Because the elasticity of the
rod was usually slightly different in different directions, the bouncing
point of light traced luminous, overlapping, and looping patterns in the
air---known today as Lissajous curves. Rayleigh references it when
discussing the compounding of two perpendicular, independent harmonic
vibrations of varying frequencies.

\textbf{Tonometer} Invented by Johann Heinrich Scheibler in 1834, the tonometer
was a highly precise acoustic apparatus used as an absolute frequency
analyzer and tuning standard before electronic measurement existed.

\textbf{How it worked}: It consisted of a large, sequentially tuned set
of precision tuning forks---often 56 or more---spanning exactly one
octave. Each fork was meticulously filed down so that it was separated
from the previous one by a tiny, identical interval (usually exactly 4
Hz), creating an audible rhythmic beating sound when played together.

\textbf{Rayleigh's Context}: By counting the exact number of beats per
second generated between an unknown acoustic pitch and the closest
tuning forks on the tonometer, scientists could mathematically calculate
the absolute frequency of any sound with extreme accuracy. Rayleigh used
these to study pure tones and verify the exact interference patterns of
wave frequencies.

Here are a few sentences guessing importance of Marin Mersenne
(1588--1648), minimizing almost everything, and otherwise exaggerating: He was
able to get a few hundred of Western Europe's best scientists to adopt
of each others work. He created a clearing house for ideas, centered
around Parisien scientists, and created or standardized Peer Review. In
acoustics, he was among the first to experimentally measure the absolute
speed of sound in air. Mersenne's Laws (1636) mathematically describe
the behavior of vibrating strings. He proved how a string's frequency
changes based on its length, cross-sectional diameter, tension, and
material density. This work serves as Rayleigh's foundation for building
the transverse string vibration equations.

\textbf{Sonometer} Also historically known as a monochord, the sonometer is an
ancient physics apparatus (dating back to Pythagoras) used to study the
mathematical laws of vibrating strings.

\textbf{How it worked}: It features a long, hollow wooden resonance box with one
or more metal wires stretched tightly across it. The wires are anchored
at one end, while the other end uses hanging weights or tuning pegs to
meticulously control the tension. A movable wooden \textit{bridge} can be slid
underneath the string to change its effective vibrating length without
changing its tension.

\textbf{Rayleigh's Context}: Plucking or bowing the string produces a clear tone.
Rayleigh used the sonometer to experimentally prove that the frequency
of a string is inversely proportional to its length, directly
proportional to the square root of its tension, and inversely
proportional to its mass density.

\textbf{Son Rauque}
Savart's \emph{son rauque} F{é}lix Savart's observed \emph{son rauque}
during 1831 experiments on the longitudinal vibrations of solid
cylindrical bars. The main tone of a bar is its longitudinal pitch.
Alongside its main tone the bar unexpectedly emits a rough,
lower-pitched note that is exactly one octave lower. In the 1850s Alfred
Terquem (1831--1887) explained son rauque as an acoustic coupling
effect. The deformation for a High frequency longitudinal
stretch/compression forces a transverse (bending) vibration. This
secondary bending wave acts as a sub-harmonic, vibrating at half the
frequency.

\textit{ Terquem collaborated with instrument makers like Rudolph Koenig}.
Koenig manufactured Terquem's \emph{verges de laiton} adopted by
the music education curriculum.

P.S.  This might be an AI hallucination.
I found no references to inventions of Terquem.
Alfred Terquem was an experimental physicists.
Koenig learning from Helmholtz to make H's apparatus',
such as resonators.
Then Koenig fabricated the sound flame, which was crucial.
I did not look for Savart's son rauque or \emph{verges de laiton}.


\textbf{J. Bourget} (1822--1887) fluids and solids, circular membrane vibration (1866),

section to the vibrations of membranes. He explicitly quotes the
collaborative experimental work of Bourget and Bernard, who tested how
circular membranes react when physically blasted with acoustic notes
from organ pipes.

Chladni's plates or diagrams? Mathieu?

Rudolph Koenig (1832--1901) transformed acoustics from a theoretical branch of
math into a precise experimental science.  He invented the manometric flame
apparatus which visualized sound waves using gas flames.  Also, he perfected
the construction of ultra-precise tuning forks and tonometers.  Rayleigh relied
on a manometric flame apparatus in his experiments.  Rayleigh cites Koenig's
empirical data on interferents between subtle overtones, and changes of pitch
in tuning forks near a resonant cavity,

\section{Torsional Vibration}\label{torsional-vibration}

When a uniform, circular bar or shaft experiences twisting along its
longitudinal axis, it undergoes torsional vibration.


\begin{longtable}[]{@{}cc@{}}
\tabularnewline%
quantity & symbol \\
\midrule\noalign{}
\endfirsthead%
\toprule\noalign{}
quantity & symbol \\
\midrule\noalign{}
\endhead%
\endlastfoot{}%
axial position & \(x\)  \\
 angle of twist & \(\theta (x,t)\) \\
torsional wave velocity & \(c = \sqrt{\frac{G}{\rho}}\) \\
\end{longtable}

Torsional vibrations share an the mathematical structure with
axial/longitudinal vibrations, with \(\theta\) instead of \(u\).
Assuming a circular cross-section, e.g., cross-sections remain perfectly
planar, the governing partial differential equation is a standard
one-dimensional wave equation:
\[
\ddot{\theta}=\frac{G}{\rho }\theta_{,xx}.
\]
In the same bar, the ratio between longitudinal and torsional wave speeds
\(\sqrt{2(1+\mu)}\).

\section{Thermoacoustics}\label{thermoacoustics}

Pieter Rijke discovered the Rijke Tube (Rye-kuh, 1859).
Singing flames and the Rijke tube phenomenon
are two variations of thermoacoustics (turning heat into sound).
Both rely on a physics principle known as the Rayleigh Criterion, which
states that if heat is added to a gas at the moment of highest
compression (maximum pressure), the acoustic vibration will be
amplified.

The Singing Flame (Higgins' Flame, 1777)

Setup: A small gas flame (originally hydrogen) is placed inside a larger
vertical glass or metal tube open at both ends.

The Mechanism: The heat from the flame causes air to expand and rush up
the tube. This airflow naturally triggers a weak, native acoustic
standing wave inside the pipe. The pressure changes of this wave force
the flame to flicker violently. Because the flame fluctuates in perfect
synchronization with the air pressure, it releases bursts of heat at
exactly the right moments to amplify the sound wave, causing the pipe to
loudly \textit{sing}.

The Rijke Tube (1859), Dutch physicist Pieter Rijke

Setup: Instead of an open flame inside the pipe, a fine metal wire mesh
is placed inside the lower half of a vertical tube.

The Mechanism: The mesh is heated with a burner until it is red-hot. The
burner is then removed.

The Result: Convection draws cold air into the bottom of the tube. As
this air passes up through the hot mesh, the mesh transfers heat to the
air in periodic pulses that sync with the pipe's natural resonant
frequency, creating an incredibly loud, deep, foghorn-like roar without
any active flame present.

In short: A singing flame uses an active, oscillating chemical reaction
to generate thermoacoustic sound, whereas a Rijke tube uses a static,
pre-heated metal element combined with buoyant convection to achieve the
exact same acoustic amplification.

\section{Appendix}\label{appendix}

The discovery of sub-harmonics in fluids, Faraday waves, motivated the
discovery of acoustic sub-harmonics. Amazing discussion of Faraday
waves, and related period states.

Professor Ratul Dasgupta, Introduction to Interfacial Waves, Lectures
\( 1,2, \dots 52, \ldots \)

\url{https://www.youtube.com/watch?v=VJ-Me2Q3tdg}

\url{https://www.youtube.com/watch?v=J6i4izpfEtw}

\url{https://www.youtube.com/watch?v=2BOGnDg9NX8}

dead link: \url{https://nptel.ac.in/course.html}

More vocabulary.

\begin{itemize}
\item
  aliquot: an `aliquot part' of an integer is an exact divisor or
  submultiple.
\item
  gutta-percha: thermoplastic latex sap.
\item
  gutta-percha tube: a hollow pipe used for communication by sound
  (stethoscopes until 1950)
\item
  diapson: sound made by a pipe organ when a chord is held.
\item
  Faulhorn: a mountain in the Alps
\item cateris paribus: all other things being equal
\item concord: a harmonious, stable combination of notes that sounds complete
and at rest to the ear.
\item ex hypothesi: according to the hypothesis proposed
\item Leyden : capacitor
\end{itemize}

Article 19. The sound is like the howl of a wolf, and is called the wolf
interval.

Capacitance is the shunt capacitance per unit length.
Conductance represents the imperfect insulation between wires, where a tiny
amount of current \textit{leaks} through the dielectric material.

contemporary RGLC parameters,  \( j = \sqrt{-1} \)

\begin{tabular}{cc}
\(\Omega \) & Ohm \\
\(H\) &  Henry  \\
\(F\) &  Farad  \\
\(R\) & \(\Omega/m\)  \\
\(G\) &  Leakage-Conductance \\
 &  Conductance (S/m) \\
\(L\) &  inductance \(H/m\) \\
\(C\) & permittance \(F/m\) \\
 & capacitance  \\
\(S\) & Siemen \\
\end{tabular}

I do not know what notation Rayleigh used.

Voltage Drop and Current Leakage equations,
\[
\frac{\partial}{\partial x}
\begin{bmatrix} V \\ I \end{bmatrix}
+
\frac{\partial}{\partial t}
\begin{bmatrix}
0 & L \\
C & 0
\end{bmatrix}
\begin{bmatrix} V \\ I \end{bmatrix} = -
\begin{bmatrix}
0 & R \\
V & 0
\end{bmatrix}
\begin{bmatrix} V \\ I \end{bmatrix}
\]

Impedance
\[
Z_{0}(\omega) =\sqrt{\frac{R+j\omega L}{G+j\omega C}}
\]

The L in LCR meter refers to an inductometer.

Vocabulary:
hundred weight (cwt), 112 lbs (force) =  1 hundred weight

\[\vec{x} = r( t - \sin t, 1 - \cos t)^T \tag{cycloid} \]
\[\vec{x} = ( r t - d \sin t, r - d \cos t)^T \tag{trochoid} \]
curtate \(d < r \) , or prolate \(d > r \) .

aeolian harp, determinental, determinate, extensible, isochronism (constant period),
Germansilver (Schlage keys), crispation, perigee, lycopodium powder (Chladni)

Smee cell  - a type of battery  circa 1900
